\begin{thebibliography}{99}
	
	% FOUNDATIONAL WORKS
	\bibitem[McCulloch \& Pitts(1943)]{McCulloch1943}
	McCulloch, W. S., \& Pitts, W. (1943). 
	A logical calculus of the ideas immanent in nervous activity. 
	\emph{The Bulletin of Mathematical Biophysics}, 5(4), 115--133.[1, 2]
	\doi{10.1007/BF02478259}.
	
	\bibitem{Rosenblatt1958}
	Rosenblatt, F. (1958). 
	The Perceptron: A probabilistic model for information storage and organization in the brain. 
	\emph{Psychological Review}, 65(6), 386--408.[3]
	\doi{10.1037/h0042519}.
	
	\bibitem[Minsky \& Papert(1969)]{Minsky1969}
	Minsky, M., \& Papert, S. (1969). 
	\emph{Perceptrons: An Introduction to Computational Geometry}. 
	MIT Press.
	
	\bibitem[Zaslavsky(1975)]{Zaslavsky1975}
	Zaslavsky, T. (1975). 
	Facing up to arrangements: face-count formulas for partitions of space by hyperplanes. 
	\emph{Memoirs of the American Mathematical Society}, 1(1), 1--102.[4]
	\doi{10.1090/memo/0154}.
	
	\bibitem[Polyak(1964)]{Polyak1964}
	Polyak, B. T. (1964).
	Some methods of speeding up the convergence of iteration methods.
	\emph{USSR Computational Mathematics and Mathematical Physics}, 4(5), 1--17.[5]
	
	\bibitem[Nesterov(1983)]{Nesterov1983}
	Nesterov, Y. (1983).
	A method for unconstrained convex minimization problem with the rate of convergence O(1/k\textasciicircum2).
	\emph{Dokl. Akad. Nauk. USSR}, 269(3), 543--547.[6]
	
	\bibitem{Rumelhart1986}
	Rumelhart, D. E., Hinton, G. E., \& Williams, R. J. (1986). 
	Learning representations by back-propagating errors. 
	\emph{Nature}, 323(6088), 533--536.[7]
	\doi{10.1038/323533a0}.
	
	\bibitem[Cybenko(1989)]{Cybenko1989}
	Cybenko, G. (1989). 
	Approximation by superpositions of a sigmoidal function. 
	\emph{Mathematics of Control, Signals and Systems}, 2(4), 303--314.[8]
	\doi{10.1007/BF02551274}.
	
	\bibitem[Vapnik(1998)]{Vapnik1998}
	Vapnik, V. (1998).
	\emph{Statistical Learning Theory}.
	Wiley-Interscience.
	
	\bibitem[Haykin(1999)]{Haykin1999}
	Haykin, S. (1999). 
	\emph{Neural Networks: A Comprehensive Foundation} (2nd ed.). 
	Prentice Hall.[9]
	
	\bibitem{Bishop2006}
	Bishop, C. M. (2006). 
	\emph{Pattern Recognition and Machine Learning}. 
	Springer.[10]
	
	\bibitem[Murphy(2012)]{Murphy2012}
	Murphy, K. P. (2012).
	\emph{Machine Learning: A Probabilistic Perspective}.
	The MIT Press.
	
	\bibitem{Shalev-Shwartz2014}
	Shalev-Shwartz, S., \& Ben-David, S. (2014).
	\emph{Understanding Machine Learning: From Theory to Algorithms}.
	Cambridge University Press.
	
	\bibitem[Goodfellow et al.(2016)]{Goodfellow2016}
	Goodfellow, I., Bengio, Y., \& Courville, A. (2016). 
	\emph{Deep Learning}. 
	MIT Press.[11]
	\doi{10.5555/3086952}.
	
	\bibitem[Nair \& Hinton(2010)]{Nair2010}
	Nair, V., \& Hinton, G. E. (2010). 
	Rectified Linear Units improve Restricted Boltzmann Machines. 
	In \emph{Proceedings of the 27th International Conference on Machine Learning (ICML-10)}, 807--814.
	
	\bibitem{Ioffe2015}
	Ioffe, S., \& Szegedy, C. (2015). 
	Batch normalization: Accelerating deep network training by reducing internal covariate shift. 
	In \emph{International Conference on Machine Learning} (pp. 448--456). 
	PMLR.
	
	\bibitem[He et al.(2015)]{He2015}
	He, K., Zhang, X., Ren, S., \& Sun, J. (2015). 
	Delving deep into rectifiers: Surpassing human-level performance on ImageNet classification. 
	In \emph{Proceedings of the IEEE International Conference on Computer Vision} (pp. 1026--1034).[12]
	\doi{10.1109/ICCV.2015.123}.
	
	\bibitem[He et al.(2016)]{He2016}
	He, K., Zhang, X., Ren, S., \& Sun, J. (2016). 
	Deep residual learning for image recognition. 
	In \emph{Proceedings of the IEEE Conference on Computer Vision and Pattern Recognition} (pp. 770--778). 
	\doi{10.1109/CVPR.2016.90}.
	
	\bibitem[Waseem et al. (2017)]{Waseem2017}
	Waseem, Z., \& Zenghui, W. (2017).
	Deep Convolutional Neural Networks for Image Classification: A Comprehensive Review
	In \emph{Neural Computation}, 29(9), 2352-2449.
	\doi{10.1162/neco_a_00990}.
	
	\bibitem{Eldan2016}
	Eldan, R., \& Shamir, O. (2016). 
	The power of depth for feedforward neural networks. 
	In \emph{Conference on Learning Theory} (pp. 907--940).[13] 
	PMLR.
	
	\bibitem{Duchi2011}
	Duchi, J., Hazan, E., \& Singer, Y. (2011).
	Adaptive subgradient methods for online learning and stochastic optimization.
	\emph{Journal of Machine Learning Research}, 12 (Article 61), 2121--2159.[14]
	
	\bibitem[Hinton(2012)]{Hinton2012}
	Hinton, G. (2012).
	Neural Networks for Machine Learning, Lecture 6.5: RMSProp.
	Coursera.
	
	\bibitem{kingma2015}
	Kingma, D. P., \& Ba, J. (2015).
	Adam: A method for stochastic optimization.
	In \emph{International Conference on Learning Representations (ICLR)}.[15]
	
	\bibitem[Ge et al.(2015)]{Ge2015}
	Ge, R., Huang, F., Jin, C., \& Yuan, Y. (2015).
	Escaping from saddle points: online stochastic gradient for tensor decomposition.
	In \emph{Conference on Learning Theory} (pp. 797--842).[16]
	
	\bibitem[Keskar et al.(2017)]{Keskar2017}
	Keskar, N. S., Mudigere, D., Nocedal, J., Smelyanskiy, M., \& Tang, P. T. P. (2017).
	On large-batch training for deep learning: Generalization gap and sharp minima.
	In \emph{International Conference on Learning Representations (ICLR)}.[17]
	
	\bibitem[Loshchilov \& Hutter(2017)]{Loshchilov2017}
	Loshchilov, I., \& Hutter, F. (2017).
	SGDR: Stochastic gradient descent with warm restarts.
	In \emph{International Conference on Learning Representations (ICLR)}.[18]
	
	\bibitem{Wilson2017}
	Wilson, A. C., Roelofs, R., Stern, M., Srebro, N., \& Recht, B. (2017).
	The marginal value of adaptive gradient methods in machine learning.
	In \emph{Advances in Neural Information Processing Systems (NIPS)} (pp. 4148--4158).
	
	\bibitem[Goyal et al.(2017)]{Goyal2017}
	Goyal, P., Dollár, P., Girshick, R., Noordhuis, P., Wesolowski, L., Kyrola, A., Tulloch, A., Jia, Y., \& He, K. (2017).
	Accurate, large minibatch SGD: Training ImageNet in 1 hour.
	\emph{arXiv preprint arXiv:1706.02677}.[19, 20]
	
	\bibitem{Wolpert1997}
	Wolpert, D. H., \& Macready, W. G. (1997). 
	No free lunch theorems for optimization. 
	\emph{IEEE Transactions on Evolutionary Computation}, 1(1), 67--82.[21]
	\doi{10.1109/4235.585893}.
	
	\bibitem{Bottou2018}
	Bottou, L., Curtis, F. E., \& Nocedal, J. (2018).
	Optimization methods for large-scale machine learning.
	\emph{SIAM Review}, 60(2), 223--311.[22]
	\doi{10.1137/16M1080173}.
	
	\bibitem{Glorot2010}
	Glorot, X., \& Bengio, Y. (2010).
	Understanding the difficulty of training deep feedforward neural networks.
	In \emph{Proceedings of the Thirteenth International Conference on Artificial Intelligence and Statistics (AISTATS)} (pp.~249--256).[23]
	PMLR.
	
	\bibitem{Reddi2018}
	Reddi, S. J., Kale, S., \& Kumar, S. (2018).
	On the convergence of Adam and beyond.
	In \emph{International Conference on Learning Representations (ICLR)}.
	
	\bibitem[Loshchilov \& Hutter(2019)]{Loshchilov2019}
	Loshchilov, I., \& Hutter, F. (2019).
	Decoupled weight decay regularization.
	In \emph{International Conference on Learning Representations (ICLR)}.[24]
	
\end{thebibliography}
